\documentclass[11pt,a4paper]{article}
\usepackage[margin=0.85in]{geometry}
\usepackage[T1]{fontenc}
\usepackage{lmodern}
\usepackage{amsmath,amssymb}
\usepackage{xcolor}
\usepackage{fancyhdr}
\usepackage{hyperref}
\setlength{\parskip}{0.35em}
\setlength{\parindent}{0pt}
\setlength{\headheight}{15pt}
\linespread{1.02}
\pagestyle{fancy}
\fancyhf{}
\lhead{CST444 Soft Computing - Module 2}
\rhead{PDF-Source Locked Guide}
\cfoot{\thepage}
\hypersetup{colorlinks=true,linkcolor=blue,urlcolor=blue}

\newcommand{\examtip}[1]{%
\vspace{0.2em}
\noindent\fcolorbox{black}{gray!12}{\parbox{0.97\linewidth}{\textbf{Exam Tip:} #1}}%
\vspace{0.2em}
}

\begin{document}

\section*{Source Lock Declaration}
This Module 2 guide is rewritten using only \textbf{SC - M2 -Ktunotes.in.pdf}. Question coverage comes from your OCR question papers, while answer points are sourced from that PDF and its training/testing problem sections.

\section*{Module 2 Topic Map from Questions}
\textbf{Frequently repeated concepts}
\begin{itemize}
\item Perceptron architecture, training algorithm, and testing algorithm.
\item Perceptron learning rule and weight updates with error correction.
\item Perceptron implementation of AND/OR/ANDNOT style logic.
\item Adaline model, delta rule (LMS/Widrow-Hoff), training and testing.
\item BPN architecture, stages, backpropagated error, and weight updation.
\item Perceptron convergence condition and XOR discussion.
\end{itemize}

\textbf{High-mark areas}
\begin{itemize}
\item BPN architecture + three-stage training explanation.
\item Perceptron 2-epoch style implementation questions.
\item Adaline training algorithm and delta-rule interpretation.
\item Error propagation significance in multilayer networks.
\end{itemize}

\textbf{Must-study-first order}
\begin{enumerate}
\item Perceptron network units and output idea.
\item Perceptron learning rule and stop condition.
\item Perceptron testing algorithm and linear separability link.
\item Adaline architecture and delta rule.
\item Adaline training/testing steps.
\item BPN architecture, feed-forward, backpropagation, and updates.
\end{enumerate}

\section*{Core Theory (From Source PDF)}

\subsection*{1) Perceptron Network Basics}
\begin{itemize}
\item Source notes describe three units: sensory, associator, and response unit.
\item Perceptron output is written as activation over net input.
\item Learning rule updates weights between associator and response unit.
\item Error is based on desired target versus calculated output.
\item Learning rate controls amount of update each step.
\item Goal is class membership decision for input patterns.
\end{itemize}

\subsection*{2) Perceptron Training and Testing}
\textbf{Training (source points)}
\begin{itemize}
\item Present patterns one by one; one full pass is one epoch.
\item Initial weights/bias are commonly set to zero in source examples.
\item Update occurs only when predicted output differs from target.
\item Stop when no weight changes in current epoch or max epoch reached.
\end{itemize}

\textbf{Testing (source points)}
\begin{itemize}
\item Testing checks performance of trained network.
\item Compute output for each input vector using learned weights.
\item Source links testing with linear separability decision boundary.
\item No training update is done in testing mode.
\end{itemize}

\subsection*{3) Adaline (Adaptive Linear Neuron)}
\begin{itemize}
\item Developed by Widrow and Hoff (source note statement).
\item Difference from perceptron: Adaline uses linear activation relation.
\item Adaline uses bipolar input and bipolar target style in source sections.
\item Trained using delta rule (also called LMS/Widrow-Hoff rule).
\item Objective is minimizing mean squared error.
\item Gradient-descent idea is explicitly discussed in notes.
\end{itemize}

\subsection*{4) Delta Rule and Adaline Algorithms}
\begin{itemize}
\item Delta rule reduces difference between net input and target value.
\item Major aim is minimizing error over all patterns.
\item Source provides stepwise training algorithm and testing algorithm.
\item Training includes initialization, net calculation, update, stopping check.
\item Testing uses trained weights, net calculation, and activation application.
\item Adaline problem sets include OR and ANDNOT training tasks.
\end{itemize}

\subsection*{5) Back Propagation Network (BPN)}
\begin{itemize}
\item BPN is multilayer feedforward network with differentiable activation.
\item Algorithm changes weights to classify input-output pairs correctly.
\item Source states training has three stages.
\item Stage 1: feed-forward of input pattern.
\item Stage 2: error calculation and backpropagation.
\item Stage 3: weight updation.
\end{itemize}

\subsection*{6) Error Propagation and Convergence View}
\begin{itemize}
\item Error is directly available at output layer.
\item Hidden-layer error is not directly available.
\item Backpropagation computes hidden corrections to reduce output error.
\item Source contrasts memorization and generalization behavior.
\item Perceptron is linked to linear separability conditions.
\item XOR is discussed in the module as a nonlinearly separable challenge context.
\end{itemize}

\examtip{For full marks, always name the three BPN stages exactly in order: feed-forward, error backpropagation, weight updation.}

\section*{Numericals and Derivations (Source-Only Procedure)}

\subsection*{N1) Perceptron OR (binary inputs, bipolar targets) -- source problem pattern}
\textbf{Given in source problem set}
\begin{itemize}
\item OR implementation is shown up to multiple epochs.
\item Initial weights/bias and learning rate setup are provided in example flow.
\end{itemize}

\textbf{Source-based steps}
\begin{itemize}
\item Present each pattern one by one.
\item Compute net input and predicted output.
\item Compare with target and check if mismatch exists.
\item Update weights only when mismatch occurs.
\item Complete all patterns to finish one epoch.
\item Repeat epoch until stop condition is satisfied.
\end{itemize}

\subsection*{N2) Adaline training (single output) -- source method}
\textbf{Given:} bipolar training pairs and initialized parameters.

\textbf{Find:} updated weights and bias after required epoch(s).

\textbf{Source-based steps}
\begin{itemize}
\item Set weights/bias and learning rate.
\item Calculate net input for each training pair.
\item Compute error factor from target and model response.
\item Adjust weights and bias according to delta-rule step.
\item Track highest weight change for stopping check.
\item Continue until tolerance condition or epoch limit.
\end{itemize}

\subsection*{N3) BPN error-propagation type question -- source method}
\textbf{Given:} multilayer architecture and target-output mismatch.

\textbf{Source-based method}
\begin{itemize}
\item Perform feed-forward to obtain output.
\item Compute output-layer error.
\item Propagate this error backward to hidden layer.
\item Use propagated corrections for weight update.
\item Repeat for all patterns and epochs.
\end{itemize}

\section*{Solved Questions for Module 2}

\subsection*{Part A / 3-Mark Questions (Each Answer Has 6 Scoring Points)}

\textbf{Q1. Explain training algorithm of perceptron network.}
\begin{enumerate}
\item Initialize weights, bias, and learning rate.
\item Present one training pattern at a time.
\item Compute net input and output using activation rule.
\item Compare calculated output with target output.
\item If error occurs, adjust trainable weights.
\item Repeat pattern loop across epochs until stop condition.
\end{enumerate}

\textbf{Q2. State testing algorithm used in perceptron network.}
\begin{enumerate}
\item Testing evaluates trained network performance.
\item Use weights obtained from training phase.
\item Apply each test input to the network.
\item Compute net input and corresponding output class.
\item Compare predicted class with expected class.
\item Report classification result or accuracy.
\end{enumerate}

\textbf{Q3. List stages involved in backpropagation network.}
\begin{enumerate}
\item Feed-forward of training input.
\item Calculation of output error.
\item Back-propagation of error to hidden layer.
\item Weight updation based on propagated error.
\item Repeat for all patterns.
\item Continue epochs until stopping criterion.
\end{enumerate}

\textbf{Q4. Explain architecture and delta-rule update in Adaline.}
\begin{enumerate}
\item Adaline is a single linear output unit model.
\item Inputs and one bias connection feed the output unit.
\item Weights are trainable and can be positive/negative.
\item Error is measured against target output.
\item Delta rule adjusts weights to reduce error.
\item Goal is minimization of mean squared error.
\end{enumerate}

\textbf{Q5. What activations are used in backpropagation networks?}
\begin{enumerate}
\item BPN needs continuously differentiable activation functions.
\item Source notes mention monotonic differentiable activation requirement.
\item Inputs/outputs can be binary or bipolar.
\item Differentiability is needed for gradient-based updates.
\item Hidden/output layers use activation suitable for backpropagation.
\item Activation choice affects convergence and learning stability.
\end{enumerate}

\textbf{Q6. How is error propagated in BPN? Explain phase-II idea.}
\begin{enumerate}
\item Error is directly computed at output layer.
\item Hidden-layer error is not directly measurable.
\item Error signals are propagated backward from output.
\item Backward correction estimates hidden contribution.
\item These corrections are used in weight updates.
\item This backward correction stage corresponds to phase-II error propagation.
\end{enumerate}

\subsection*{Part B / 14-Mark Questions (Each Answer Has 14 Scoring Points)}

\textbf{Q1. Explain significance of error terms in BPN and explain architecture.}
\begin{enumerate}
\item BPN is a multilayer feedforward network.
\item It contains input, hidden, and output layers.
\item Hidden and output layers include bias terms.
\item Training objective is reducing target-output mismatch.
\item Output-layer error is directly computed.
\item Hidden-layer error is inferred by backward propagation.
\item Source states this is required for hidden-layer weight tuning.
\item Stage 1 is feed-forward computation.
\item Stage 2 is error calculation and backpropagation.
\item Stage 3 is weight updation.
\item This process repeats for all training patterns.
\item Repeated epochs reduce overall output error.
\item Differentiable activation is required in BPN.
\item Conclusion: error terms guide all layer-wise weight corrections.
\end{enumerate}

\textbf{Q2. Implement OR with binary inputs and bipolar targets using perceptron (up to epochs).}
\begin{enumerate}
\item Use source OR-perceptron problem setup.
\item Initialize weights/bias as shown in source pattern.
\item Set learning rate as prescribed.
\item Present one input-target pair at a time.
\item Compute output from perceptron rule.
\item Check mismatch with bipolar target.
\item Update trainable weights when mismatch occurs.
\item Update bias similarly.
\item Continue for all four OR patterns (one epoch).
\item Repeat for next epoch(s).
\item Stop when no further updates are needed.
\item Verify all OR patterns are classified correctly.
\item Report final learned weights and bias from table/process.
\item Conclusion: OR converges under perceptron due separability.
\end{enumerate}

\textbf{Q3. Explain Adaline training algorithm for single-output class.}
\begin{enumerate}
\item Initialize random non-zero weights and bias.
\item Set learning-rate parameter in source-recommended range.
\item Present each bipolar training pair.
\item Compute net input to output unit.
\item Compute error with respect to target.
\item Apply delta-rule correction to each weight.
\item Update bias along with weights.
\item Repeat for all patterns in one epoch.
\item Monitor highest weight change in epoch.
\item Compare with tolerance condition.
\item Stop when change is below tolerance.
\item Otherwise continue epochs.
\item Use final trained parameters for testing.
\item Conclusion: Adaline learns by continuous error minimization.
\end{enumerate}

\textbf{Q4. Implement one epoch of Adaline for logic function (source style).}
\begin{enumerate}
\item Write given input-target table (binary/bipolar as question states).
\item Record initial $w$, $b$, and learning-rate values.
\item For first pattern, compute net response.
\item Compute error factor against target.
\item Update each weight using delta-rule correction.
\item Update bias using same correction principle.
\item Move to second pattern and repeat.
\item Continue similarly for all patterns in epoch.
\item Maintain tabular record of each update step.
\item At epoch end, report updated parameters.
\item Mention whether stop condition is met.
\item If not met, indicate next epoch continuation.
\item Keep calculations in sequence to secure method marks.
\item Conclusion: one-epoch result is an intermediate trained state.
\end{enumerate}

\textbf{Q5. Implement AND with binary inputs and bipolar targets using perceptron.}
\begin{enumerate}
\item Use source AND-perceptron worked-problem structure.
\item Set initial weights and bias.
\item Write bipolar target mapping for AND.
\item Compute output for each pattern.
\item Compare with target for each row.
\item Update weights for mismatch rows.
\item Update bias accordingly.
\item Complete first epoch updates.
\item Repeat epoch if needed.
\item Stop when all outputs match targets.
\item Verify final classification for four input rows.
\item State final weights and bias from process table.
\item Mention separability reason for convergence.
\item Conclusion: perceptron correctly learns AND boundary.
\end{enumerate}

\textbf{Q6. Draw BPN architecture and explain training algorithm.}
\begin{enumerate}
\item Draw input, hidden, and output layers.
\item Show forward signal direction for feed-forward phase.
\item Show reverse direction for error backpropagation phase.
\item Initialize all trainable weights and biases.
\item Feed one input vector and compute outputs.
\item Compare output with target and compute error.
\item Backpropagate error to hidden layer.
\item Update output-layer incoming weights.
\item Update hidden-layer incoming weights.
\item Repeat for all training patterns.
\item Complete one epoch and evaluate progress.
\item Continue until stopping criterion.
\item Note requirement of differentiable activation function.
\item Conclusion: BPN training is iterative gradient-based correction.
\end{enumerate}

\textbf{Q7. What is Adaline? Draw model and explain OR/ANDNOT training approach.}
\begin{enumerate}
\item Adaline means adaptive linear neuron.
\item Source describes it as single linear unit with trainable weights.
\item Bias acts as an adjustable weight input.
\item Training uses delta rule (LMS/Widrow-Hoff).
\item Write bipolar input-target table for required gate.
\item Compute net response for each pattern.
\item Compute error term from target-response mismatch.
\item Apply weight and bias updates pattern-wise.
\item Continue for all patterns in epoch.
\item Repeat for required number of epochs.
\item Evaluate progress using error reduction idea.
\item Use final weights for testing.
\item Verify gate behavior on all input combinations.
\item Conclusion: Adaline improves learning smoothness versus threshold-only update.
\end{enumerate}

\textbf{Q8. Explain perceptron convergence idea and discuss XOR.}
\begin{enumerate}
\item Source connects perceptron with linear separability concept.
\item Training update continues when errors occur.
\item Stopping can occur when no updates are needed in epoch.
\item This behavior corresponds to separable-pattern convergence.
\item XOR is used in module as nonlinearly separable example context.
\item Non-separable patterns cannot be separated by single linear boundary.
\item Therefore simple perceptron is insufficient for XOR-type separation.
\item Multilayer methods are introduced to address this gap.
\item BPN is one such multilayer method in this module.
\item BPN uses backpropagated error for hidden-layer adjustment.
\item Differentiable units enable gradient correction.
\item This allows non-linear decision mapping.
\item Exam sketch of XOR point layout strengthens justification.
\item Conclusion: perceptron convergence is separability-dependent; XOR motivates multilayer learning.
\end{enumerate}

\section*{Formula/Method Sheet (Source-Safe)}
\begin{itemize}
\item Perceptron output is activation over net input.
\item Perceptron update occurs only when calculated output differs from target.
\item Learning rate controls step size in updates.
\item Epoch: one complete presentation of all patterns.
\item Adaline training uses delta rule (LMS/Widrow-Hoff).
\item Delta-rule objective: minimize mean squared error.
\item BPN stages: feed-forward, error backpropagation, weight updation.
\item BPN requires monotonic differentiable activation.
\end{itemize}

\section*{One-Page Quick Revision}
\begin{itemize}
\item Perceptron: error-triggered weight correction for separable classes.
\item Testing phase never changes trained weights.
\item Adaline: linear unit + delta rule + error minimization.
\item BPN: hidden-layer correction through backward error flow.
\item Learn these three phrases exactly: epoch, learning rate, stop condition.
\item For long answers, always state architecture first, then algorithm steps.
\item For numerical questions, show row-wise update table clearly.
\item Mention linear-separability link when discussing perceptron behavior.
\end{itemize}

\examtip{When you are short on time, write algorithm steps as numbered lines first; this secures method marks before arithmetic detail.}

\end{document}
